% ============================================================
%  Panel Review Report
%  Document: Stochastic Degradation Models for Train Fleet
%            Maintenance Scheduling
%  Date: 2026-05-13
% ============================================================
\documentclass[11pt,a4paper]{article}

% ---- packages ----
\usepackage[utf8]{inputenc}
\usepackage[T1]{fontenc}
\usepackage{lmodern}
\usepackage{amsmath,amssymb}
\usepackage[margin=2.2cm]{geometry}
\usepackage{booktabs}
\usepackage{enumitem}
\usepackage{xcolor}
\usepackage{hyperref}
\usepackage{longtable}
\usepackage{tabularx}
\usepackage{array}
\usepackage{multicol}
\usepackage{fancyhdr}
\usepackage{titlesec}
\usepackage{tocloft}

% ---- colours ----
\definecolor{acceptgreen}{RGB}{0,100,0}
\definecolor{judgmentorange}{RGB}{180,95,6}
\definecolor{resolvedblue}{RGB}{30,80,160}
\definecolor{lightgray}{RGB}{245,245,245}
\definecolor{sectionblue}{RGB}{20,60,120}

% ---- header/footer ----
\pagestyle{fancy}
\fancyhf{}
\fancyhead[L]{\small\textcolor{gray}{Panel Review Report}}
\fancyhead[R]{\small\textcolor{gray}{2026-05-13}}
\fancyfoot[C]{\small\thepage}
\renewcommand{\headrulewidth}{0.4pt}

% ---- section formatting ----
\titleformat{\section}{\Large\bfseries\color{sectionblue}}{}{0em}{}[\titlerule]
\titleformat{\subsection}{\large\bfseries\color{sectionblue!80}}{}{0em}{}

% ---- custom commands ----
\newcommand{\acceptlabel}{\textbf{\textcolor{acceptgreen}{[ACCEPT]}}}
\newcommand{\judgmentlabel}{\textbf{\textcolor{judgmentorange}{[JUDGMENT]}}}
\newcommand{\resolvedlabel}{\textbf{\textcolor{resolvedblue}{[RESOLVED]}}}
\newcommand{\commentbox}[1]{\noindent\colorbox{lightgray}{\parbox{\dimexpr\textwidth-2\fboxsep}{#1}}}

% ---- metadata ----
\title{\vspace{-1.5cm}\textbf{Formal Panel Review Report}\\[6pt]
\large Stochastic Degradation Models for Train Fleet Maintenance Scheduling:\\
Taxonomy, Literature, and Mathematical Foundations}
\author{Review Panel: Mathematician $\cdot$ Domain Expert $\cdot$ Enthusiast $\cdot$ Devil's Advocate $\cdot$ Computer Scientist}
\date{May 13, 2026}

\begin{document}
\maketitle
\thispagestyle{fancy}

\tableofcontents
\newpage

% ============================================================
\section{Executive Summary}
% ============================================================

\subsection{Document Assessed}
The reviewed document is a comprehensive literature review paper (${\sim}2{,}181$ lines) that constructs a dual taxonomy (macro and micro) of stochastic orders in optimization, then positions a train fleet maintenance scheduling MILP with FSD loop constraints within these taxonomies. It surveys six literature streams and claims genuine novelty for the combination of fleet-level scheduling, parametric stochastic degradation (Gamma, Wiener, IG), and stochastic-dominance sustainability constraints.

\subsection{Panel Composition}
Five specialist reviewers provided independent critiques:
\begin{itemize}[nosep,leftmargin=1.5em]
  \item \textbf{Mathematician} --- Mathematical correctness, proof rigor, notation consistency (20 comments)
  \item \textbf{Domain Expert} --- Operations research, reliability engineering, practical relevance (18 comments)
  \item \textbf{Enthusiast} --- Research opportunities, extensions, connections (12 comments)
  \item \textbf{Devil's Advocate} --- Overclaims, logical gaps, missing literature, scope limitations (25 comments)
  \item \textbf{Computer Scientist} --- Complexity, scalability, MILP formulation quality, algorithms (23 comments)
\end{itemize}

\subsection{Summary Statistics}

\begin{table}[h]
\centering
\begin{tabular}{lcc}
\toprule
\textbf{Metric} & \textbf{Count} \\
\midrule
Total comments across all reviewers & 88 \\
Group A (Must Accept) & 27 \\
Group B (Requires Judgment) & 37 \\
Debate clusters resolved & 8 (all consensus in round 1) \\
\midrule
\multicolumn{2}{l}{\textit{Group A by severity}} \\
\quad Major & 7 \\
\quad Moderate & 10 \\
\quad Minor & 10 \\
\midrule
\multicolumn{2}{l}{\textit{Group B by severity}} \\
\quad Major & 4 \\
\quad Moderate & 18 \\
\quad Minor & 15 \\
\bottomrule
\end{tabular}
\caption{Review statistics overview.}
\end{table}

\subsection{High-Level Verdict}
The paper makes a genuine organizational contribution through its dual-taxonomy architecture and the positioning of a novel research direction. The mathematical content is largely sound, with definitions correctly stated and the key parametric FSD reductions for Gamma and IG families correctly argued. However, the panel identified \textbf{27 items requiring mandatory correction}, including:
\begin{itemize}[nosep,leftmargin=1.5em]
  \item A mathematical error in the supermartingale-FSD hierarchy (Remark~1.1)
  \item An internal inconsistency in the Wiener FSD claim (the MILP enforces increasing convex order, not FSD)
  \item An unverified assumption underlying the IG tractability claim
  \item An incorrect binary variable count
  \item Multiple missing derivations in the MILP formulation
  \item Missing critical literature references
\end{itemize}
The panel recommends \textbf{major revision}. No theorem is outright false, but several claims are imprecise, insufficiently justified, or internally inconsistent. All 8 disagreements were resolved by consensus in a single debate round.

\newpage
% ============================================================
\section{Panel Comments --- Must Accept (Group A)}
% ============================================================

Items in this section meet at least one of: mathematical error, undefined term, missing assumption, logical gap, factual contradiction, internal inconsistency, notation ambiguity, missing citation for non-trivial claim, structural defect, or unanimous panel agreement.

% ---- A-01 ----
\subsection{A-01: Supermartingale-FSD Hierarchy Error}
\acceptlabel\quad\textit{Severity: moderate}\quad\textit{Location: Remark 1.1, lines 139--153}

\textbf{Issue.} Remark~1.1 claims a three-level hierarchy: supermartingale $\Rightarrow$ marginal FSD $\Rightarrow$ mean non-increase. The implication supermartingale $\Rightarrow$ marginal FSD is \emph{mathematically incorrect}.

A supermartingale satisfies $\mathbb{E}[X_{(k+1)T}\mid\mathcal{F}_{kT}]\le X_{kT}$ a.s., which implies $\mathbb{E}[X_{(k+1)T}]\le\mathbb{E}[X_{kT}]$ but does \emph{not} imply $\mathcal{L}(X_{(k+1)T})\le_{\mathrm{st}}\mathcal{L}(X_{kT})$. The Mathematician provides a counterexample: $X_0=1$ a.s., $X_1=0$ with probability $2/3$ and $X_1=3$ with probability $1/3$ is a martingale where the CDFs cross.

\textbf{Proposed modification.} Replace the linear chain with: ``Both marginal FSD and the supermartingale property imply mean non-increase, but neither implies the other in general. They are two distinct strengthenings along different axes (distributional vs.\ conditional).''

\textit{Reviewer: Mathematician (M-5). Verified: yes.}

\bigskip

% ---- A-02 ----
\subsection{A-02: Wiener FSD Claim Is Self-Contradicted}
\acceptlabel\quad\textit{Severity: major}\quad\textit{Location: Section 3.2, Wiener bullet, lines 1533--1548}

\textbf{Issue.} The paper acknowledges that strict FSD between Gaussians with unequal variances is impossible (CDFs cross), yet claims FSD ``reduces to finite parameter conditions'' for the Wiener model. The two-constraint formulation ($\mu_{2H}\le\mu_H$ and $v_{2H}\le v_H$) does \emph{not} enforce FSD when $v_{2H}<v_H$. The Mathematician provides a numerical counterexample confirming the CDFs cross.

\textbf{Debate resolution (Cluster~5, consensus).} The Enthusiast identified that the Wiener constraints are necessary and sufficient for the \emph{increasing convex order} ($\le_{\mathrm{icx}}$) between Gaussian marginals --- a well-defined stochastic ordering strictly weaker than FSD (Shaked and Shanthikumar 2007, Thm~4.A.22).

\textbf{Proposed modification.}
\begin{enumerate}[nosep,leftmargin=1.5em]
  \item Add remark in Section~3.2: ``For Gaussian marginals, the conditions $\mu_{2H}\le\mu_H$ and $v_{2H}\le v_H$ are necessary and sufficient for the increasing convex order ($\le_{\mathrm{icx}}$), strictly weaker than FSD.''
  \item Update Table~5: replace ``Parametric FSD reduction'' with ``Parametric SD reduction: FSD (Gamma, IG) / increasing convex order (Wiener).''
  \item Add citation to Shaked and Shanthikumar (2007), Theorem~4.A.22.
  \item Note that $\le_{\mathrm{icx}}$ implies non-increase of expected cost under any increasing convex loss function but does not prevent certain tail events from worsening.
\end{enumerate}

\textit{Reviewers: Mathematician (M-3), Devil's Advocate (DA-02), Enthusiast (MC4). Verified: yes.}

\bigskip

% ---- A-03 ----
\subsection{A-03: Lemma 3.1 Should Be Relabeled as Remark}
\acceptlabel\quad\textit{Severity: moderate}\quad\textit{Location: Lemma 3.1, lines 1756--1768}

\textbf{Issue.} Lemma~3.1 (Wiener variance monotonicity under ARD$_1$) follows directly from the variance dynamics definition in Eq.~(6): imperfect repair leaves $v$ unchanged, missions add positive increments, so $v$ is non-decreasing. This is a modeling consequence, not a mathematical result about Wiener processes.

\textbf{Debate resolution (Cluster~1, consensus).} All four debating reviewers agreed: relabel as Remark, remove novelty claims, preserve practical discussion motivating the IG alternative.

\textbf{Proposed modification.}
\begin{enumerate}[nosep,leftmargin=1.5em]
  \item Rename ``Lemma~3.1'' to ``Remark.''
  \item Add: ``This follows directly from the variance dynamics in Eq.~(6), which define imperfect repair as leaving $v$ unchanged.''
  \item Remove the phrase ``novel structural result'' if used in reference to this remark.
  \item Retain the discussion of practical implications for MILP design and the IG workaround.
\end{enumerate}

\textit{Reviewers: Mathematician (M-2), Devil's Advocate (DA-16), Enthusiast (S4), Domain Expert (S2). Verified: yes.}

\bigskip

% ---- A-04 ----
\subsection{A-04: Binary Variable Count Is Incorrect}
\acceptlabel\quad\textit{Severity: minor}\quad\textit{Location: Section 3.4, Challenge 3, lines 1786--1794}

\textbf{Issue.} The formula $F(M+L)\cdot 2H$ omits the $+1$ for the maintenance activity ($j=0$) and the factor of~2 for separate repair vs.\ replacement variables. Three reviewers independently computed the correct formula:
\[
\text{Binary variables} = F\cdot 2H\cdot(M+1+2L).
\]
For $F=20$, $M=15$, $L=5$, $H=180$: $20\times 360\times 26 = 187{,}200$, not the stated $144{,}000$.

\textbf{Proposed modification.} Correct the formula and numerical example. Add a complete variable/constraint census table.

\textit{Reviewers: Mathematician (M-17), Computer Scientist (CS-M1), Domain Expert (M5). Verified: yes.}

\bigskip

% ---- A-05 ----
\subsection{A-05: IG FSD Reduction Requires Unverified Fixed-Ratio Assumption}
\acceptlabel\quad\textit{Severity: major}\quad\textit{Location: Section 3.2, IG bullet, lines 1518--1528}

\textbf{Issue.} The reduction of IG FSD to a single linear inequality $\mu_{2H}\le\mu_H$ is conditioned on $\lambda/\mu$ being held constant. The paper acknowledges this caveat but:
\begin{itemize}[nosep,leftmargin=1.5em]
  \item Does not verify which parameterization the MILP uses (fixed $\lambda$ vs.\ fixed $\lambda/\mu$).
  \item If $\lambda$ is fixed per component (the natural physical case), the reduction is incorrect.
  \item The reference to Ye and Chen (2014) ``Sec.~4'' is too vague --- that paper is about parameter estimation, not stochastic ordering.
\end{itemize}

\textbf{Proposed modification.} Clarify the IG parameterization. Provide either a self-contained proof or a specific theorem citation. If the MILP uses fixed $\lambda$, derive the correct (bivariate) FSD condition.

\textit{Reviewers: Mathematician (M-1), Domain Expert (W4), Devil's Advocate (DA-03), Computer Scientist (CS-J2). Verified: consensus on the gap.}

\bigskip

% ---- A-06 ----
\subsection{A-06: Lattice Completeness Conflates Tail Probability with Support}
\acceptlabel\quad\textit{Severity: moderate}\quad\textit{Location: Proposition 1.4, lines 388--411}

\textbf{Issue.} The claim that the degradation threshold $\tau$ provides ``compact support $[0,\tau]$'' is incorrect. The reliability constraint $\Pr(D>\tau)\le\varepsilon$ bounds the tail probability but does not truncate the support. The distribution still has support on all of $\mathbb{R}_+$.

\textbf{Debate resolution (Cluster~2, consensus).} Downgrade Proposition~1.3 to a Conjecture with three explicit hypotheses: lattice completeness, order-preservation of $\Phi_\theta$, and Galois connection for co-design.

\textbf{Proposed modification.} Replace ``compact support $[0,\tau]$'' with a correct argument using parametric subfamilies. Cite Mueller and Scarsini (2006) for the precise completeness result.

\textit{Reviewers: Domain Expert (W6), Mathematician (M-4). Verified: yes.}

\bigskip

% ---- A-07 ----
\subsection{A-07: Order-Preservation of $\Phi_\theta$ Not Established}
\acceptlabel\quad\textit{Severity: moderate}\quad\textit{Location: Section 3.2, Branch 5, lines 1573--1586}

\textbf{Issue.} The Knaster-Tarski application requires order-preservation (stochastic monotonicity) of $\Phi_\theta$, which is not proven. The Mercier and Castro (2019) citation establishes a different ordering result (policy comparison), not kernel monotonicity.

\textbf{Proposed modification.} Add a one-paragraph proof for Gamma/ARD$_1$: (a) monotone transformation by $(1-\rho)$ preserves $\le_{\mathrm{st}}$; (b) convolution with independent non-negative increment preserves $\le_{\mathrm{st}}$. Correct the Mercier-Castro citation scope.

\textit{Reviewer: Mathematician (M-10). Debate resolution: Cluster~2 (consensus).}

\bigskip

% ---- A-08 ----
\subsection{A-08: Big-M Linearization Never Made Explicit}
\acceptlabel\quad\textit{Severity: major}\quad\textit{Location: Lines 1425--1436}

\textbf{Issue.} The big-M linearization of bilinear terms ($\mu_{il,k-1}\cdot x^m_{ilk}$) is mentioned but never made explicit. The paper never states: (a)~what $M$ is, (b)~how many auxiliary variables/constraints each linearization introduces, (c)~whether the big-M values are tight. LP relaxation quality depends entirely on big-M tightness.

\textbf{Proposed modification.} Add the standard big-M reformulation with four linear inequalities. Derive $M$ from reliability constraints (e.g., $M=\tau$ for Wiener mean). State tightness.

\textit{Reviewer: Computer Scientist (CS-M2). Verified: yes.}

\bigskip

% ---- A-09 ----
\subsection{A-09: Three-Case Dynamics Linearization Never Shown}
\acceptlabel\quad\textit{Severity: major}\quad\textit{Location: Lines 1406--1409}

\textbf{Issue.} The three-case dynamics (no action / imperfect repair / full replacement) requires modeling three mutually exclusive cases with indicator constraints or additional big-M constraints. The full linearization is never shown.

\textbf{Proposed modification.} Write out the full linearized dynamics for one time step, showing all auxiliary variables and big-M constraints (${\sim}4$ constraints and 1 auxiliary per bilinear term per $(i,l,k)$ triple).

\textit{Reviewer: Computer Scientist (CS-M6). Verified: yes.}

\bigskip

% ---- A-10 ----
\subsection{A-10: Wiener Reliability Linear Approximation Never Derived}
\acceptlabel\quad\textit{Severity: major}\quad\textit{Location: Lines 1345--1361}

\textbf{Issue.} The linear approximation $(\Phi^{-1}(1-\varepsilon))^2 v + 2\tau\mu\le\tau^2$ of the SOC reliability constraint is stated without derivation. Squaring is not generally a conservative operation.

\textbf{Proposed modification.} Derive the approximation from the rotated SOC constraint, or cite Ben-Tal and Nemirovski (2001) with a specific theorem number (not the entire book).

\textit{Reviewers: Computer Scientist (CS-M5), Domain Expert (M6). Verified: yes.}

\bigskip

% ---- A-11 ----
\subsection{A-11: Missing Rolling Stock Maintenance Routing Literature}
\acceptlabel\quad\textit{Severity: major}\quad\textit{Location: Section 4.3, lines 1945--1961}

\textbf{Issue.} The novelty claim about routing-maintenance co-optimization ignores the established rolling stock maintenance routing literature: Maroti and Kroon (2005, Transportation Science), Wagenaar et al.\ (2017, Transportation Science). These jointly optimize vehicle routing and maintenance scheduling.

\textbf{Debate resolution (Cluster~6, consensus).} Add references and a paragraph discussing how FSD loop constraints differ from threshold-based triggers in that literature.

\textbf{Proposed modification.} Add Maroti and Kroon (2005), Wagenaar et al.\ (2017) to Section~4.3. Add clarifying paragraph.

\textit{Reviewers: Domain Expert (W1), Devil's Advocate (DA-06). Verified: yes.}

\bigskip

% ---- A-12 through A-27: shorter entries ----

\subsection{A-12: Damage Regularizer $u$: $C_D>0$ and Range Not Stated}
\acceptlabel\quad\textit{Severity: minor}\quad\textit{Location: Definition 3.2, lines 1395--1423}

The epigraph reformulation $u=\max_{i,k}\sum_l\mu_{ilk}$ is exact only when $C_D>0$. The range of $(i,k)$ is unspecified. State $C_D>0$ as a requirement and specify the index range.

\textit{Reviewers: Mathematician (M-16), Domain Expert (M4).}

\subsection{A-13: Citation Key / Year Mismatches}
\acceptlabel\quad\textit{Severity: minor}\quad\textit{Location: references.bib}

Four citation keys have wrong years: \texttt{olde2014}$=$2017, \texttt{dariano2017}$=$2019, \texttt{uithetbroek2019}$=$2021, \texttt{andradeGomes2015}$=$2012. Harmonize all keys with actual publication years.

\textit{Reviewers: Mathematician (M-19), Domain Expert (W9), Devil's Advocate (DA-19).}

\subsection{A-14: Indexing Convention Not Documented}
\acceptlabel\quad\textit{Severity: minor}\quad\textit{Location: Definition 3.1, lines 1259--1379}

Activity index $j$ is 0-based ($j=0$ for maintenance) while all other indices are 1-based. Add a notational convention paragraph at the start of Section~3.1.

\textit{Reviewer: Mathematician (M-15).}

\subsection{A-15: Periodic Extension Boundary Error at $k=H$}
\acceptlabel\quad\textit{Severity: minor}\quad\textit{Location: Lines 1309--1310}

At $k=H$, $k\bmod H=0$, which is outside the defined range $[1,H]$. Use $((k-1)\bmod H)+1$ or equivalent.

\textit{Reviewer: Domain Expert (M2).}

\subsection{A-16: Tightness Argument Incomplete}
\acceptlabel\quad\textit{Severity: moderate}\quad\textit{Location: Section 3.3, Pillar 4, lines 1625--1640}

The family $\{\mathcal{L}(D_{i,l,kH})\}_{k\ge1}$ is tight because all members are stochastically bounded above by $\mathrm{Gamma}(\hat\alpha,\beta)$ and below by $\delta_0$. Cite Prokhorov's theorem explicitly.

\textit{Reviewer: Mathematician (M-8).}

\subsection{A-17: Common-Scale Requirement for Gamma FSD}
\acceptlabel\quad\textit{Severity: moderate}\quad\textit{Location: Section 3.2, Gamma bullet, lines 1511--1515}

The Gamma FSD reduction to $\alpha_1\le\alpha_2$ requires common scale $\beta$. This assumption is stated but buried. Box or bold the assumption and note that violating it invalidates the single-constraint reduction.

\textit{Reviewer: Mathematician (M-9).}

\subsection{A-18: Blanket Claim About Weak LP Relaxations}
\acceptlabel\quad\textit{Severity: major}\quad\textit{Location: Lines 1791--1794}

The assignment subproblem is a transportation polytope (TU), giving a strong LP relaxation. Weakness comes only from the dynamics linearization. Replace the blanket claim with a nuanced statement.

\textit{Reviewer: Computer Scientist (CS-M3).}

\subsection{A-19: Mission Coverage Constraints Not Stated}
\acceptlabel\quad\textit{Severity: minor}\quad\textit{Location: Line 1269, Definition 3.1}

The constraint $\sum_i x_{ijk}=1$ for each mission $j\ge1$ and each $k$ is implied but never formally stated. Add it to the formulation.

\textit{Reviewer: Domain Expert (M3).}

\subsection{A-20: Book-Level References Too Vague}
\acceptlabel\quad\textit{Severity: minor}\quad\textit{Location: Lines 1428--1435}

Cite Conforti et al.\ (2014) \S10.1 and Ben-Tal and Nemirovski (2001) \S3.2 for the specific linearization and SOCP claims.

\textit{Reviewer: Computer Scientist (CS-S3).}

\subsection{A-21: FSD CDF Direction Convention}
\acceptlabel\quad\textit{Severity: moderate}\quad\textit{Location: Definition 1.2, lines 169--179}

The convention $F_X(t)\ge F_Y(t)$ for $X\le_{\mathrm{st}}Y$ is correct but opposite to many readers' expectation. Add a clarifying sentence in the degradation context.

\textit{Reviewers: Mathematician (M-6), Domain Expert (M8).}

\subsection{A-22: SSD/Increasing Convex Order Sign Convention}
\acceptlabel\quad\textit{Severity: minor}\quad\textit{Location: Definition 1.4, lines 233--248}

Clarify: ``$Y$ SSD-dominates $X$'' in finance corresponds to $X\le_{\mathrm{icx}}Y$.

\textit{Reviewers: Mathematician (M-20), Domain Expert (M1).}

\subsection{A-23: Implication Chain Omits Reversed Hazard-Rate Order}
\acceptlabel\quad\textit{Severity: minor}\quad\textit{Location: Proposition 1.1, lines 297--315}

Add that $\le_{\mathrm{lr}}\Rightarrow\le_{\mathrm{rh}}$ and forward-reference Figure~2 for the complete hierarchy.

\textit{Reviewer: Mathematician (M-13).}

\subsection{A-24: Gaussian Ordering Alternatives Conflated}
\acceptlabel\quad\textit{Severity: minor}\quad\textit{Location: Remark 2.1, lines 747--766}

Distinguish: (a) $\le_{\mathrm{icx}}$ requires $\mu_1\le\mu_2$ and $\mu_1^2+\sigma_1^2\le\mu_2^2+\sigma_2^2$; (b) $\le_{\mathrm{cx}}$ requires equal means and $\sigma_1\le\sigma_2$.

\textit{Reviewer: Mathematician (M-14).}

\subsection{A-25: Single-Period Sufficiency Needs Formal Proof}
\acceptlabel\quad\textit{Severity: minor}\quad\textit{Location: Lines 1373--1378}

Add a one-paragraph induction proof with three explicit assumptions: identical policy repetition, stochastic monotonicity, exact initial condition transfer.

\textit{Reviewer: Mathematician (M-18).}

\subsection{A-26: FSD Testing Complexity Imprecise}
\acceptlabel\quad\textit{Severity: minor}\quad\textit{Location: Lines 1155--1160}

Replace ``$O(n)$'' with ``$O(n)$ on common ordered support; $O(n\log n)$ for distinct supports.''

\textit{Reviewer: Computer Scientist (CS-S6).}

\subsection{A-27: Wiener Reliability Approximation (Domain Expert)}
\acceptlabel\quad\textit{Severity: minor}\quad\textit{Location: Lines 1429--1433}

Provide brief derivation or specific Ben-Tal and Nemirovski (2001) theorem citation for the linear relaxation. (Overlaps with A-10.)

\textit{Reviewer: Domain Expert (M6).}


\newpage
% ============================================================
\section{Panel Comments --- Requires Your Judgment (Group B)}
% ============================================================

Items in this section involve scope expansion, stylistic preferences, extensions, boundary conditions where scope is debatable, or trade-off recommendations. Each entry identifies who raised it, any dissenting views, and the rationale for requiring judgment.

% ---- B-01 ----
\subsection{B-01: Novelty Claim Framing}
\judgmentlabel\quad\textit{Severity: major}\quad\textit{Location: Section 5, Table 5, Table 6}

\textbf{Issue.} The Devil's Advocate argues the novelty claim is tautological: the taxonomy axes are tailored so only the target application checks every box. The Domain Expert finds the claim ``likely correct but overstated.'' The Enthusiast finds it convincingly demonstrated.

\textbf{Debate resolution (Cluster~3, consensus).} Narrow the claim: acknowledge rolling stock maintenance routing literature; specify novelty is in FSD loop constraints, not routing-maintenance coupling per se.

\textbf{Trade-off.} Narrowing the claim is more honest but reduces perceived impact. The user must decide how aggressively to reframe.

\textit{Raised by: Devil's Advocate (DA-01, DA-10). Dissent: Enthusiast (S5).}

\bigskip

% ---- B-02 ----
\subsection{B-02: Four Contributions vs.\ One Structured Survey}
\judgmentlabel\quad\textit{Severity: major}\quad\textit{Location: Abstract, lines 73--82}

\textbf{Issue.} The Devil's Advocate argues all four ``contributions'' are components of a single literature review. The Enthusiast implicitly accepts the multi-contribution framing.

\textbf{Debate resolution (Cluster~7, consensus).} Reframe as one primary contribution comprising four components.

\textbf{Trade-off.} This is an editorial framing choice. The content is not contested.

\textit{Raised by: Devil's Advocate (DA-22). Dissent: Enthusiast.}

\bigskip

% ---- B-03 ----
\subsection{B-03: Lattice-Theoretic Framing Scope}
\judgmentlabel\quad\textit{Severity: high}\quad\textit{Location: Branch 5; Propositions 1.3--1.4}

\textbf{Issue.} The Enthusiast values the lattice-theoretic connection; the Devil's Advocate calls it decorative; the Mathematician, CS, and Domain Expert identify fixable gaps (completeness, order-preservation, Galois connection, non-constructivity).

\textbf{Debate resolution (Cluster~2, consensus).} Retain but qualify: downgrade to Conjecture, fix support argument, verify Gamma/ARD$_1$ case, label Censi as analogy, note non-constructivity.

\textbf{Trade-off.} Full verification requires new mathematical work. Retaining without verification (as qualified conjecture) is acceptable but less satisfying.

\textit{Raised by: all five reviewers with varying assessments.}

\bigskip

% ---- B-04 ----
\subsection{B-04: Taxonomy Scope --- Branches 4 and 5}
\judgmentlabel\quad\textit{Severity: moderate}\quad\textit{Location: Sections 1.2.4--1.2.5}

\textbf{Issue.} Branches 4 (Wasserstein) and 5 (Lattice) are not used by the fleet problem. The Devil's Advocate sees padding; the Enthusiast values reference breadth.

\textbf{Debate resolution (Cluster~4, consensus).} Retain but label: ``included for taxonomic completeness; not directly activated by the fleet problem.'' Condense speculative application paragraphs.

\textbf{Trade-off.} Breadth vs.\ focus depends on target audience.

\textit{Raised by: Devil's Advocate (DA-25), Domain Expert (W3, W10). Dissent: Enthusiast (S1).}

\bigskip

% ---- B-05 ----
\subsection{B-05: Practical Justification for FSD over Thresholds}
\judgmentlabel\quad\textit{Severity: major}\quad\textit{Location: Section 3.4, Challenge 1}

\textbf{Issue.} No concrete example shows when threshold constraints fail but FSD prevents unsustainable degradation. The Domain Expert and Devil's Advocate both flag this. No reviewer directly contests the importance of this gap.

\textbf{Trade-off.} Adding a worked example would strengthen the paper significantly but requires new content (possibly computational experiments). The scope depends on whether the paper remains a pure literature review or incorporates illustrative results.

\textit{Raised by: Domain Expert (W2), Devil's Advocate (DA-09, DA-17, DA-24).}

\bigskip

% ---- B-06 ----
\subsection{B-06: Schedule Periodicity and Sustainability Guarantee}
\judgmentlabel\quad\textit{Severity: moderate}\quad\textit{Location: Definition 3.1}

\textbf{Issue.} The MILP does not enforce schedule periodicity. Nothing prevents completely different structures in $[1,H]$ and $[H{+}1,2H]$. The sustainability guarantee requires policy repetition, which is not enforced. In rolling-horizon implementation, the guarantee vanishes.

\textbf{Trade-off.} Adding periodicity constraints changes the formulation. Providing a formal argument without them requires new analysis. Acknowledging the limitation honestly may be the simplest path.

\textit{Raised by: Domain Expert (W7), Devil's Advocate (DA-07, DA-08, DA-13, DA-15), Mathematician (M-18).}

\bigskip

% ---- B-07 ----
\subsection{B-07: Independence Assumption and Dependence Material}
\judgmentlabel\quad\textit{Severity: moderate}\quad\textit{Location: Section 2.5, Section 3.3 Pillar 5}

\textbf{Debate resolution (Cluster~8, consensus).} Condense from ${\sim}80$ to ${\sim}15$--$20$ lines. Promote independence to named Assumption. Sketch copula extension as future direction.

\textbf{Trade-off.} Condensation amount is an editorial judgment.

\textit{Raised by: Devil's Advocate (DA-04), Domain Expert (W8), Enthusiast (O5).}

\bigskip

% ---- B-08 through B-37: condensed entries ----

\subsection{B-08: Literature Survey Systematic Exclusions}
\judgmentlabel\quad\textit{Severity: moderate}

Omits DRO-CBM, chance-constrained fleet maintenance, nonlinear Wiener processes, model selection. Debate resolution (Cluster~6): add targeted paragraphs. \textit{Raised by: Devil's Advocate (DA-06, DA-14), Domain Expert (W5), Enthusiast (O1).}

\subsection{B-09: Overreach in Application Connections}
\judgmentlabel\quad\textit{Severity: moderate}

Commodity pricing, MCMC, mean-field games connections are hypothetical. Condense to brief mentions. \textit{Raised by: Domain Expert (W3), Devil's Advocate (DA-12, DA-18).}

\subsection{B-10: Decomposition Opportunity Not Developed}
\judgmentlabel\quad\textit{Severity: major}

No Benders/DW/Lagrangian decomposition sketched despite block-diagonal structure. Requires significant new algorithmic content for a literature review. \textit{Raised by: Computer Scientist (CS-M4), Enthusiast (O7).}

\subsection{B-11: Supermodularity Verification}
\judgmentlabel\quad\textit{Severity: moderate}

Either verify or disprove increasing differences. CS notes this is straightforward. \textit{Raised by: Devil's Advocate (DA-11), CS (CS-J3), Enthusiast (O6).}

\subsection{B-12: IG Process $t^2$ Scaling}
\judgmentlabel\quad\textit{Severity: moderate}

The Mathematician flags that for i.i.d.\ IG increments, convolution gives $\mathrm{IG}(t\eta,t\lambda)$, not $\mathrm{IG}(t\eta,t^2\lambda)$. \textbf{Unverified} --- requires author verification. \textit{Raised by: Mathematician (M-11).}

\subsection{B-13: Sato (1999) Theorem Number}
\judgmentlabel\quad\textit{Severity: minor}

Levy-Ito decomposition may be Thm~19.3, not 19.2. Unverified; depends on edition. \textit{Raised by: Mathematician (M-12).}

\subsection{B-14--B-37: Additional Items}
The remaining Group~B items (24 entries) cover:
\begin{itemize}[nosep,leftmargin=1.5em]
  \item \textbf{B-14}: Deterministic parameter framing (stylistic, Devil's Advocate)
  \item \textbf{B-15}: Damage regularizer computational effect (CS suggestion)
  \item \textbf{B-16}: Paper scope clarification: review vs.\ specification (CS)
  \item \textbf{B-17}: Example 1.1 HR/FSD distinction (Mathematician)
  \item \textbf{B-18}: Explicit convergence rate bounds (Enthusiast extension)
  \item \textbf{B-19}: Wiener FSD relaxation hierarchy (Enthusiast extension)
  \item \textbf{B-20}: NP-hardness and H-sweep elaboration (CS)
  \item \textbf{B-21}: Wasserstein-DRO connection (CS)
  \item \textbf{B-22}: Cutting-plane complexity characterization (CS)
  \item \textbf{B-23}: Symmetry-breaking from periodic extension (CS)
  \item \textbf{B-24}: Assignment polytope TU implications (CS)
  \item \textbf{B-25}: Objective coefficient sensitivity (CS)
  \item \textbf{B-26}: Conclusion lacks new computational insight (CS)
  \item \textbf{B-27}: Maintenance capacity and loop feasibility (Devil's Advocate)
  \item \textbf{B-28}: Wiener/Gaussian terminology inconsistency (Domain Expert)
  \item \textbf{B-29--B-30}: Figure 1 density and Table 6 font size (Domain Expert)
  \item \textbf{B-31}: Section 4.6 title (Domain Expert)
  \item \textbf{B-32--B-36}: Questions for authors (Domain Expert: infeasibility handling, Wiener-IG validity, computational testing, Galois connection, empirical process theory)
  \item \textbf{B-37}: Tarski non-constructivity (CS; addressed by Cluster~2 resolution)
\end{itemize}


\newpage
% ============================================================
\section{Debate Resolutions}
% ============================================================

All eight disagreements identified during integration were resolved by \textbf{consensus in a single round} of structured debate. No second round was needed.

\begin{longtable}{p{1.2cm}p{4cm}p{3.5cm}p{5cm}}
\toprule
\textbf{Cluster} & \textbf{Topic} & \textbf{Outcome} & \textbf{Key Resolution} \\
\midrule
\endhead
1 & Lemma 3.1 significance & Consensus (1 round) & Relabel as Remark. Remove novelty claims. Preserve practical discussion. \\
\addlinespace
2 & Lattice/Knaster-Tarski framing & Consensus (1 round) & Downgrade Proposition~1.3 to Conjecture. Fix support argument. Verify Gamma/ARD$_1$ case. Label Censi as analogy. \\
\addlinespace
3 & Novelty claim validity & Consensus (1 round) & Narrow to FSD-specific innovation. Acknowledge rolling stock routing literature. Reframe Table~6 as literature mapping. \\
\addlinespace
4 & Taxonomy scope (Branches 4--5) & Consensus (1 round, residual concern) & Retain but label for completeness. Condense speculative paragraphs. Add DRO connection. DA residual: overall length reduction. \\
\addlinespace
5 & Wiener two-constraint formulation & Consensus (1 round) & Identified as increasing convex order ($\le_{\mathrm{icx}}$), not FSD. Update Table~5 and add Shaked-Shanthikumar citation. \\
\addlinespace
6 & Literature survey completeness & Consensus (1 round) & Add rolling stock maintenance routing references. Add DRO/chance-constrained paragraph. Add Ye and Xie (2015). \\
\addlinespace
7 & Contribution count & Consensus (1 round) & One primary contribution with four components. Add scope clarification. \\
\addlinespace
8 & Dependence orders material & Consensus (1 round) & Condense to 15--20 lines. Promote independence to named Assumption. Sketch copula extension. Move formal definitions to appendix. \\
\bottomrule
\caption{Summary of all debate resolutions.}
\end{longtable}

\textbf{Convergence pattern.} Most disagreements stemmed from a tension between the paper's theoretical ambition and taxonomic breadth (valued by the Enthusiast) and the gap between that ambition and rigorous execution or practical relevance (flagged by the Mathematician, Devil's Advocate, and Domain Expert). A recurring resolution pattern was to retain the broad framing but add explicit qualifications, verify formal conditions, and narrow overclaims. The Enthusiast's observations frequently provided constructive paths through disagreements (e.g., identifying the Wiener constraint as increasing convex order, proposing copula extensions).

\newpage
% ============================================================
\section{Proposed Modifications}
% ============================================================

This section consolidates all proposed text and content changes. Items from Group~A are labeled \acceptlabel{} (mandatory). Items from Group~B are labeled \judgmentlabel{} (pending user decision).

\subsection{Mathematical Corrections}

\begin{enumerate}[label=\arabic*.,leftmargin=2em]
  \item \acceptlabel{} \textbf{Remark 1.1}: Fix supermartingale-FSD hierarchy. Replace linear chain with branching structure. (A-01)
  \item \acceptlabel{} \textbf{Section 3.2, Wiener}: Identify constraint as $\le_{\mathrm{icx}}$, not FSD. Update Table~5. (A-02)
  \item \acceptlabel{} \textbf{Lemma 3.1}: Relabel as Remark. (A-03)
  \item \acceptlabel{} \textbf{Challenge 3}: Correct variable count to $187{,}200$. (A-04)
  \item \acceptlabel{} \textbf{Section 3.2, IG}: Verify or correct fixed-ratio assumption. (A-05)
  \item \acceptlabel{} \textbf{Proposition 1.4}: Fix support/tail-probability conflation. (A-06)
  \item \acceptlabel{} \textbf{Branch 5}: Prove order-preservation for Gamma/ARD$_1$. (A-07)
  \item \acceptlabel{} \textbf{Remark 2.1}: Distinguish $\le_{\mathrm{icx}}$ from $\le_{\mathrm{cx}}$ for Gaussians. (A-24)
  \item \acceptlabel{} \textbf{Lines 1373--1378}: Add formal induction proof for propagation. (A-25)
  \item \acceptlabel{} \textbf{Lines 1309--1310}: Fix boundary case $k\bmod H=0$. (A-15)
  \item \judgmentlabel{} \textbf{Example 2.3}: Verify IG $t^2$ scaling vs.\ convolution formula. (B-12)
  \item \judgmentlabel{} \textbf{Proposition 2.5}: Verify Sato (1999) theorem number. (B-13)
\end{enumerate}

\subsection{MILP Formulation Completeness}

\begin{enumerate}[label=\arabic*.,leftmargin=2em]
  \item \acceptlabel{} Show big-M linearization explicitly with bounds derived from reliability constraints. (A-08)
  \item \acceptlabel{} Show full three-case dynamics linearization. (A-09)
  \item \acceptlabel{} Derive Wiener reliability linear approximation. (A-10, A-27)
  \item \acceptlabel{} Add mission coverage constraints to formulation. (A-19)
  \item \acceptlabel{} State $C_D>0$ requirement; specify index range for $u$. (A-12)
  \item \acceptlabel{} Replace blanket LP relaxation weakness claim. (A-18)
  \item \judgmentlabel{} Add decomposition sketch (Benders/Lagrangian). (B-10)
  \item \judgmentlabel{} Add scope clarification: review vs.\ specification. (B-16)
\end{enumerate}

\subsection{Notation and Conventions}

\begin{enumerate}[label=\arabic*.,leftmargin=2em]
  \item \acceptlabel{} Add indexing convention paragraph. (A-14)
  \item \acceptlabel{} Clarify FSD CDF direction in degradation context. (A-21)
  \item \acceptlabel{} Clarify SSD/finance sign convention. (A-22)
  \item \acceptlabel{} Add $\le_{\mathrm{rh}}$ to Proposition 1.1. (A-23)
  \item \acceptlabel{} Fix FSD testing complexity to $O(n\log n)$. (A-26)
  \item \acceptlabel{} Specify book sections for references. (A-20)
  \item \acceptlabel{} Fix bibliography year mismatches. (A-13)
  \item \judgmentlabel{} Adopt consistent Wiener/Gaussian terminology. (B-28)
\end{enumerate}

\subsection{Framing and Scope}

\begin{enumerate}[label=\arabic*.,leftmargin=2em]
  \item \acceptlabel{} \textbf{Section 3.3, Pillar 4}: Complete tightness argument with Prokhorov. (A-16)
  \item \acceptlabel{} \textbf{Gamma FSD}: Make common-scale $\beta$ assumption prominent. (A-17)
  \item \acceptlabel{} Add rolling stock maintenance routing references. (A-11)
  \item \judgmentlabel{} Narrow novelty claim to FSD-specific innovation. (B-01)
  \item \judgmentlabel{} Reframe ``four contributions'' as one contribution, four components. (B-02)
  \item \judgmentlabel{} Qualify lattice-theoretic framing extensively. (B-03)
  \item \judgmentlabel{} Label Branches 4--5 for taxonomic completeness. (B-04)
  \item \judgmentlabel{} Add practical FSD vs.\ threshold example. (B-05)
  \item \judgmentlabel{} Address periodicity/sustainability gap. (B-06)
  \item \judgmentlabel{} Condense dependence material; promote independence assumption. (B-07)
  \item \judgmentlabel{} Add DRO, chance-constrained literature. (B-08)
  \item \judgmentlabel{} Condense speculative application connections. (B-09)
\end{enumerate}


\newpage
% ============================================================
\appendix
\section{Individual Reviewer Summaries}
% ============================================================

\subsection{Mathematician}
\begin{itemize}[nosep,leftmargin=1.5em]
  \item \textbf{Verdict:} Major revision. No theorem outright false, but ${\sim}8$--$10$ points need tighter formal hygiene.
  \item \textbf{Confidence:} 0.88
  \item \textbf{Comments:} 20 (3 major, 7 moderate, 10 minor)
  \item \textbf{Top priorities:} IG FSD verification (M-1), supermartingale-FSD error (M-5), Wiener FSD characterization (M-3), IG $t^2$ scaling (M-11), Lemma~3.1 reframing (M-2).
  \item \textbf{Assessment:} Largely sound mathematical content with imprecise or insufficiently justified claims on ${\sim}8$--$10$ specific points. The parametric FSD reductions for Gamma are correct; the IG and Wiener reductions need attention.
\end{itemize}

\subsection{Domain Expert}
\begin{itemize}[nosep,leftmargin=1.5em]
  \item \textbf{Verdict:} Major revision. Score: 6/10.
  \item \textbf{Confidence:} High
  \item \textbf{Comments:} 18 (4 major, 6 moderate, 8 minor) + 5 questions for authors
  \item \textbf{Top priorities:} Missing rolling stock routing literature (W1), no practical FSD justification (W2), tenuous application connections (W3), IG tractability gap (W4).
  \item \textbf{Assessment:} Taxonomic contribution is solid and positioning well-executed. Novelty claim likely correct but insufficiently defended on practical grounds. The paper reads more as theoretical positioning than practitioner-useful literature review.
\end{itemize}

\subsection{Enthusiast}
\begin{itemize}[nosep,leftmargin=1.5em]
  \item \textbf{Verdict:} Strong accept. Score: 9/10.
  \item \textbf{Confidence:} High
  \item \textbf{Comments:} 12 (6 strengths, 12 opportunities, 4 missing connections)
  \item \textbf{Top extensions:} DR-FSD branch (O1), Wiener relaxation hierarchy (O2), convergence rate bounds (O4), Benders decomposition (O7), predictive maintenance with Bayesian updates (O10).
  \item \textbf{Assessment:} Remarkably well-structured and ambitious. The dual-taxonomy architecture is genuinely original. The parametric FSD reduction is the central computational insight. Opens numerous exciting research avenues.
\end{itemize}

\subsection{Devil's Advocate}
\begin{itemize}[nosep,leftmargin=1.5em]
  \item \textbf{Verdict:} Major revision required due to overclaims.
  \item \textbf{Confidence:} High
  \item \textbf{Comments:} 25 (7 major, 13 moderate, 5 minor)
  \item \textbf{Top invalidation risks:} Wiener FSD self-contradiction (DA-02), tautological novelty claim (DA-01), IG fixed-ratio assumption (DA-03).
  \item \textbf{Assessment:} Competent literature review that repeatedly overstates the strength and novelty of its positioning. The taxonomy is designed to make the application look unique. Several sections (dependence orders, Branches~4--5) are padding. The sustainability guarantee holds only under the non-operational assumption of exact policy repetition.
\end{itemize}

\subsection{Computer Scientist}
\begin{itemize}[nosep,leftmargin=1.5em]
  \item \textbf{Verdict:} Weak. Overall score: 4/10 for computational content.
  \item \textbf{Confidence:} High
  \item \textbf{Comments:} 23 (6 must-fix, 12 should-fix, 4 judgment calls)
  \item \textbf{Scores:} Complexity analysis: 3, scalability: 2, algorithmic completeness: 3, LP relaxation: 1, big-M tightness: 1, decomposition: 1, reproducibility: 3.
  \item \textbf{Assessment:} Strong on mathematical taxonomy, weak on computational content. For a paper claiming MILP tractability as a central contribution, the computational analysis is superficial. Variable count is wrong. Big-M never derived. No decomposition algorithm sketched. No LP relaxation analysis.
\end{itemize}

\bigskip
\begin{center}
\rule{0.5\textwidth}{0.4pt}\\[6pt]
\textit{End of Panel Review Report}
\end{center}

\end{document}
